\section{Datos y metodologia}

\subsection{Fuentes de datos}

Se integran tres bloques:
\begin{itemize}
  \item \textbf{R02 (EPHC personas):} ingresos laborales, formalidad y caracteristicas sociodemograficas.
  \item \textbf{R01 (EPHC hogares):} conectividad, calidad material de vivienda y movilidad del hogar.
  \item \textbf{Serie macro mensual:} IPC general, IPC alimentos y fechas de ajuste del salario minimo legal.
\end{itemize}

Los indicadores de hogares se agregan por trimestre y se unen con los de R02 via etiqueta temporal \texttt{YYYYTrimQ}. Para evitar comparaciones invalidas, se diferencia explicitamente entre unidad persona (R02) y unidad hogar (R01).

\subsection{Variables de interes}

\begin{itemize}
  \item \textbf{Franja salarial relativa al SML:} \texttt{Menos de 1 SML}, \texttt{1 SML}, \texttt{Mas de 1 SML}, con tolerancia base de $\pm 10\%$.
  \item \textbf{Formalidad laboral:} proxy de cotizacion a caja/IPS.
  \item \textbf{Indicadores R01:} porcentaje de hogares con internet, material habitacional apto y movilidad propia.
  \item \textbf{Inflacion mensual:} variacion porcentual del indice respecto al mes previo para IPC general y alimentos.
\end{itemize}

\subsection{Estrategia empirica de evento}

Para cada evento de ajuste del SML se construye una ventana simetrica de $h$ meses pre y post (especificacion principal: $h=6$). El estadistico principal es:
\begin{equation}
EI_e = \overline{\pi}^{\,post}_e - \overline{\pi}^{\,pre}_e,
\end{equation}
donde $\pi$ representa inflacion mensual y $e$ denota evento.

Se reporta ademas un indicador de traspaso acumulado:
\begin{equation}
CT_e = \frac{EI_e}{\Delta SML_e},
\end{equation}
con $\Delta SML_e$ en puntos porcentuales.

La inferencia descriptiva se complementa con placebo Monte Carlo: se sortean ventanas ``no-evento'' de igual tamano y se compara el $EI$ observado frente a la distribucion simulada.

\subsection{Configuracion principal}

La configuracion central del paper usa eventos no solapados y $h=6$, con \nEventosMain{} eventos. Los resultados clave se resumen en la Tabla~\ref{tab:key}.

\begin{table}[H]
\centering
\caption{Resultados clave de la especificacion principal}
\label{tab:key}
\begin{tabular}{lr}
\toprule
Indicador & Valor\\
\midrule
Eventos identificados (total) & 23.000\\
Eventos usados (h=6, no solapados) & 22.000\\
EI IPC general (pp) & -0.067\\
EI IPC alimentos (pp) & -0.054\\
CT general (EI/ajuste) & -0.003\\
CT alimentos (EI/ajuste) & 0.023\\
p-value placebo IPC general & 0.581\\
p-value placebo IPC alimentos & 0.835\\
\bottomrule
\end{tabular}

\end{table}
